\documentclass[11pt]{amsart}
\usepackage[T1]{fontenc}
\usepackage{lmodern}
\usepackage{amsmath,amssymb,mathtools}
\usepackage{aliascnt}
\usepackage{microtype}
\usepackage[margin=1.15in]{geometry}
\usepackage[colorlinks=true,linkcolor=blue,citecolor=blue,urlcolor=blue]{hyperref}
\urlstyle{same}
\newcommand{\doi}[1]{\href{https://doi.org/#1}{\nolinkurl{doi:#1}}}
\usepackage[nameinlink,noabbrev]{cleveref}
\hypersetup{pdftitle={Integral tail signatures and the Hilbert--Smith conjecture},
  pdfauthor={A. Dabrowski's AI agents},
  pdfsubject={Integral signature obstruction to effective p-adic manifold actions},
  pdfkeywords={Hilbert-Smith conjecture, controlled algebra, equivariant signature}}

\numberwithin{equation}{section}
\newtheorem{theorem}{Theorem}[section]
\newaliascnt{proposition}{theorem}
\newtheorem{proposition}[proposition]{Proposition}
\aliascntresetthe{proposition}
\newaliascnt{lemma}{theorem}
\newtheorem{lemma}[lemma]{Lemma}
\aliascntresetthe{lemma}
\newaliascnt{corollary}{theorem}
\newtheorem{corollary}[corollary]{Corollary}
\aliascntresetthe{corollary}
\theoremstyle{definition}
\newaliascnt{definition}{theorem}
\newtheorem{definition}[definition]{Definition}
\aliascntresetthe{definition}
\theoremstyle{remark}
\newaliascnt{remark}{theorem}
\newtheorem{remark}[remark]{Remark}
\aliascntresetthe{remark}
\DeclareMathOperator{\Fix}{Fix}
\DeclareMathOperator{\prop}{prop}
\DeclareMathOperator{\Kar}{Kar}
\DeclareMathOperator{\Free}{Free}
\DeclareMathOperator{\Cone}{Cone}
\DeclareMathOperator{\Sign}{Sign}
\DeclareMathOperator{\Sym}{Sym}
\DeclareMathOperator{\im}{im}
\DeclareMathOperator{\tr}{tr}
\newcommand{\Q}{\mathbb Q}
\newcommand{\Z}{\mathbb Z}
\newcommand{\C}{\mathbb C}
\newcommand{\R}{\mathbb R}
\newcommand{\CP}{\mathbb {CP}}
\newcommand{\Zp}{\mathbb Z_p}
\newcommand{\Linf}{L^{\langle-\infty\rangle}}
\newcommand{\Lspec}{\mathbf L^{\langle-\infty\rangle}}

\title[Integral tail signatures and Hilbert--Smith]{Integral tail signatures and the Hilbert--Smith conjecture}
% Author information: name, affiliation, and contact address.
\author{A. Dabrowski's AI agents}
\email{alex@dabrow.cc}
\date{October 8, 2026}

\begin{document}
\begin{abstract}
The obstruction to an effective $p$-adic manifold action is an integer
signature tail. On asymptotically controlled rational modules over one
fixed free $C_p$-space, finite-cover Mayer--Vietoris makes the difference
between $p$ copies and quotient-permutation tensoring nilpotent. Every
degree-four class therefore has eventually $p$-divisible signatures.
An effective action gives a class of signature $1$: coordinate averaging
provides degree-one sphere control, free sheets turn quotient representatives
into a compatible homotopy idempotent, and its normalized Poincar\'e corner
forgets to a scalar manifold germ. Successive localization boundaries
recover one copy of $\CP^2$, up to sign. The resulting contradiction
proves the Hilbert--Smith conjecture.
\end{abstract}
\maketitle

\section{Introduction}\label{sec:introduction}

The Hilbert--Smith conjecture extends Hilbert's fifth problem from
locally Euclidean groups to groups acting on manifolds. The cases of
one- and two-manifolds are classical \cite{Pa19}; Pardon proved the
three-dimensional case \cite{Pa13}. In arbitrary dimension, the
conjecture holds for Lipschitz actions \cite{RS97} and quasiconformal
actions \cite{Mar99}. We prove the following.

\begin{theorem}[$p$-adic exclusion]\label{thm:main}
For every prime $p$, a continuous action of $\Zp$ on a connected
finite-dimensional manifold has nontrivial kernel.
\end{theorem}

\begin{corollary}[Hilbert--Smith]\label{cor:HS}
A locally compact Hausdorff group acting continuously and effectively
on a connected finite-dimensional manifold is a Lie group with its
given topology.
\end{corollary}

Manifolds are topological and Hausdorff, and may have boundary.
Lie groups are allowed arbitrarily many connected components. The
construction takes place in one coordinate chart; it requires no global
orientation or triangulation of the acted-on manifold.
The standard reduction to $p$-adic exclusion uses locally compact group
structure and Newman's theorem \cite{Pa19}.

The obstruction is an integer signature tail. For cyclic groups
$G_i=C_{p^{a_i}}$ with index-$p$ subgroups $P_i$, and a fixed compact
free $C_p$-space $T$, tensoring with the positive permutation form
$\tau_i=\Q[G_i/P_i]$ is locally $p$ copies of the identity.
Controlled Mayer--Vietoris therefore makes $p\,\mathrm{id}-\tau\otimes(-)$
nilpotent, and equivariant signature characters force every class in
$L_4(\mathcal A_{\mathbf G}(T))$ to have eventually $p$-divisible
ordinary signatures.

An effective $\Zp$-action, however, would supply a class with signature $1$.
Coordinate averaging gives an invariant degree-one sphere control map.
With this control, fine local chain models turn finite quotient
representatives into a homotopy action. Adding free sheets cancels the
detector displacement, allowing averaging to produce a duality-compatible
homotopy idempotent. Its normalized corner forgets to the scalar manifold
germ of $S^n\times\CP^2$. Successive localization boundaries then recover
one copy of $\CP^2$, up to sign. The resulting signature contradicts the
divisibility above and proves the theorem.

The construction and boundary evaluation use Karoubi localization
\cite{CP95}, controlled Poincar\'e--Lefschetz duality \cite{Ran99},
and algebraic Poincar\'e pairs and unions \cite{Ran80I,Ran80II}.
Idempotent splitting uses \cite{BS01}; the local carrier and normalization
calculations are given below.

\subsection*{Conventions}
All actions are jointly continuous. Chain coefficients are rational,
but the $L$-groups and signature targets are not rationalized.
Complexes have a single finite degree range, independent of the sequence
index; ranks and coefficient sizes are unrestricted.

In the asymptotic categories below, equality means eventual equality of
matrices, so finite initial segments of sequences may be discarded.
In an exterior quotient, equality means that the defect factors through
an exterior-supported object by asymptotically small maps. We use maximum
product metrics and take empty suprema to be zero.

We use ultimate lower quadratic $L$-theory,
$L_k(\mathcal A)=\pi_k\Lspec(\mathcal A)$
\cite[Definition 4.16]{CP95}, \cite[\S17]{Ran92}.
The properties used below are functoriality and additivity,
Karoubi localization with natural connecting maps
\cite[Theorem 4.2]{CP95}, projective-to-ultimate decoration maps,
and invariance under idempotent completion
\cite[Theorem 17.2 and Definition 17.7]{Ran92}.
For localization boundaries we use the pair description in
\cite[proof of Theorem 4.1]{CP95}. Small skeletons (one representative per
isomorphism class) avoid set-theoretic size issues in $K$- and $L$-theory.

\section{Asymptotic categories and support}\label{alg:categories}

Fix $G_i=C_{p^{a_i}}$, $a_i\ge1$, index-$p$ subgroups $P_i$, and
$G_i/P_i=C_p$. Let $X$ be a compact metric $C_p$-space with isometric
action; each $G_i$ acts through $G_i/P_i$.
\begin{definition}\label{alg:category}
An object of $\mathcal A_{\mathbf G}(X)$ is a sequence of finite free
based modules $M_i=\Q[B_i]$, with $B_i$ a free $G_i$-set and equivariant
labels $\ell_i:B_i\to X$. Morphisms are sequences of equivariant rational
matrices, modulo eventual equality, satisfying
\begin{equation}\label{alg:propagation}
 \prop(u_i)=\sup_{(u_i)_{b'b}\ne0}d(\ell_i'(b'),\ell_i(b))\longrightarrow0.
\end{equation}
Duality fixes labeled bases and transposes matrices. Trivial groups give
$\mathcal A_{\mathbf1}(X)$. Ranks and coefficients are unrestricted.
\end{definition}

The propagation bounds
\begin{equation}\label{alg:propbounds}
 \prop(vu)\le\prop(v)+\prop(u),\qquad
 \prop(u+v)\le\max\{\prop(u),\prop(v)\}
\end{equation}
show that composition and addition are well defined. Forgetting
equivariance and pushing labels through uniformly continuous equivariant
maps give duality-preserving functors.
For a closed invariant $S\subset X$, let
$\mathcal A_S=\mathcal A_{\mathbf G,S}(X)$ consist of objects
whose labels approach $S$ uniformly, with $\mathcal A_\varnothing=0$.
Write $I_S$ for the ideal of maps factoring through $\mathcal A_S$.
For a closed invariant exterior $Z\subset X$, put
\begin{equation}\label{alg:exterior}
 \mathcal B_{\mathbf G,Z}(X)=\mathcal A_{\mathbf G}(X)/\mathcal A_Z.
\end{equation}
A map belongs to $I_S$ exactly when both endpoints of every nonzero
entry approach $S$ uniformly: necessity follows from propagation, and
sufficiency by restricting the source to its nonzero columns.

\begin{lemma}[Support calculus]\label{alg:filtration}\label{alg:excision}
\label{alg:relabel}\label{alg:separation}
The inclusion $\mathcal A_S\subset\mathcal A_{\mathbf G}(X)$ is a
duality-invariant Karoubi filtration
\cite[Definition 1.27]{CP95}. For a closed invariant cover $X=A\cup B$,
with $Y=A\cap B$, there are involution-preserving equivalences
\begin{equation}\label{alg:excision-equivalence}
 \mathcal A_A\cap\mathcal A_B=\mathcal A_Y,\qquad
 \mathcal A_A/\mathcal A_Y\simeq\mathcal A_{\mathbf G}(X)/\mathcal A_B.
\end{equation}
For nonempty closed invariant $S$,
$\mathcal A_{\mathbf G}(S)\simeq\mathcal A_S$.
If $d(S,Z)>0$, $\mathcal A_S\to\mathcal B_{\mathbf G,Z}(X)$ is fully
faithful.
\end{lemma}
\begin{proof}
For a null sequence $r_i\ge0$, split each basis at distance $r_i$ from
$S$. These equivariant orthogonal splits are filtered by maxima and
compatible with sums. A map to or from an object supported within
distance $a_i$ of $S$, with propagation $\epsilon_i$, factors through
the near summand for $r=a+\epsilon$. This verifies the filtration axioms.

By compactness, points approaching both $A$ and $B$ approach $Y$.
Splitting labels into $A$ and $B\setminus A$ gives essential
surjectivity of the quotient equivalence. For a factorization between
$A$-supported objects through a $B$-supported object, retain the
intermediate rows reached by the first factor. They approach both sets,
so the factorization is through $\mathcal A_Y$. This proves full
faithfulness, including $Y=\varnothing$ by compact separation.

Relabel each free orbit using a nearest point of $S$ and extend
equivariantly. The identity matrices are natural orthogonal
isomorphisms with propagation tending to zero; relabeling changes a
map's propagation by at most its two labeling errors. Finally, the
endpoint description of $I_Z$ makes every map between $S$-supported
objects in that ideal eventually zero when $d(S,Z)>0$.
\end{proof}

Splits precede idempotent completion, denoted $\Kar$; an object $(M,e)$
carries the support of $M$.

\section{Integer signatures and tails}\label{alg:tails}

For a finite group $G$, the correspondence between invariant rational
forms $b$ and group-ring forms uses the unnormalized trace:
\begin{equation}\label{alg:trace}
 \lambda(v,w)=\sum_{g\in G}b(v,gw)g,\qquad \tr_1\lambda(v,w)=b(v,w).
\end{equation}
Thus the coefficient-of-identity trace $\tr_1$ makes forgetting
equivariance duality-preserving, including on projective summands. The half-unit comparison and \cite[Theorem 5.2]{Ran89}
identify the projectively decorated group $L_4^p(\Q[G])$ with the Witt group
of nonsingular invariant symmetric rational forms. Such a form has
signature representation
\begin{equation}\label{alg:signature}
 \Sign_G(V,b)=[V^+_{\C}]-[V^-_{\C}]\in R_{\C}(G),
\end{equation}
where $R_\C(G)$ is the complex representation ring and $V^\pm$ are the
positive and negative spectral subspaces on $V\otimes_\Q\R$ for an
invariant positive inner product. The representation is independent of
this choice: projection along one negative subspace identifies any two
positive subspaces, by inertia, and similarly for negative ones.

To see that it descends to the Witt group, let $L$ be an invariant
Lagrangian and choose an invariant complement $K$ by averaging.
Replacing $u\in K$ by $u-c(u)/2$, where $c:K\to L$ satisfies
$b(cu,v)=b(u,v)$ for $u,v\in K$, makes $K$ Lagrangian. Hence metabolic
forms have zero signature representation. The character of $\Sign_G(V,b)$
at the identity is the ordinary integer signature $\sigma(V,b)$; tensoring
with a positive form multiplies characters.

\begin{lemma}[Cyclic arithmetic]\label{alg:cyclotomic}
For $G=C_{p^a}=\langle g\rangle$, if $\Sign_G(V,b)(g)=0$, then
$p\mid\sigma(V,b)$.
\end{lemma}
\begin{proof}
Write the signature character as $\sum_{j=0}^{p^a-1}m_j\chi_j$,
$m_j\in\Z$, with $\chi_j(g)=\zeta_{p^a}^j$, where $\zeta_{p^a}$ is a
primitive $p^a$-th root of unity. Then $F(t)=\sum_jm_jt^j$ is divisible
in $\Z[t]$ by the $p^a$-th cyclotomic polynomial $\Phi_{p^a}(t)$, so
$F(1)=\sigma(V,b)$ is divisible by $\Phi_{p^a}(1)=p$.
\end{proof}

Write $\Free(R)$ for the category of finite free based $R$-modules.
Put $\mathcal P=\prod_i\Free(\Q[G_i])$ and let $\mathcal S$ be its
finite-support subcategory. Forgetting labels gives a functor to
\begin{equation}\label{alg:Vcategory}
 \mathcal V_{\mathbf G}=\mathcal P/\mathcal S,
\end{equation}
whose morphisms are sequences modulo eventual equality.

\begin{proposition}[Tail extraction]\label{alg:extraction}
There is a natural homomorphism
\begin{equation}\label{alg:tail-isomorphism}
 L_4(\mathcal V_{\mathbf G})\longrightarrow
 \prod_iL_4^p(\Q[G_i])/\bigoplus_iL_4^p(\Q[G_i]),
\end{equation}
given by the coordinate classes of any lift to $L_4(\mathcal P)$.
It commutes with forgetting equivariance and positive permutation
tensor functors.
\end{proposition}
\begin{proof}
The semisimple rings $\Q[G_i]$ have vanishing negative $K$-groups
\cite[III, \S4]{Wei13}, also as free-module categories
\cite[\S11, p.~101]{Ran92}. Bass's polynomial definition commutes with
$\mathcal S=\bigcup_{F\text{ finite}}\prod_{i\in F}\Free(\Q[G_i])$:
every finite list of objects, idempotents, polynomial matrices and
relations lies in one finite block. Consequently $K_{-k}(\mathcal S)=0$
for $k\ge1$. Rothenberg comparison then gives
$L_j(\mathcal S)=\bigoplus_iL_j^p(\Q[G_i])$ and
$L_j(\Free(\Q[G_i]))=L_j^p(\Q[G_i])$.

Splitting at finite index sets is an invariant Karoubi filtration.
Coordinate projections detect $L_3(\mathcal S)$, so its map to
$L_3(\mathcal P)$ is injective. Localization therefore lifts every
$L_4(\mathcal V_{\mathbf G})$ class to $L_4(\mathcal P)$; two lifts
differ by the image of $L_4(\mathcal S)$. Their coordinate classes
agree eventually. These constructions are natural for the stated
functors.
\end{proof}

Applying ordinary signature after extraction defines
\begin{equation}\label{alg:signature-tail}
 \sigma_{\mathrm{tail}}:L_4(\mathcal A_{\mathbf G}(X))\longrightarrow
 \Z^{\mathbb N}/\Z^{(\mathbb N)}.
\end{equation}
Here $\Z^{(\mathbb N)}$ denotes the finitely supported integer sequences.
We use this map only after localization has removed the exterior.

\section{Divisibility over a free \texorpdfstring{$C_p$}{C(p)}-space}
\label{alg:obstruction}

Let $T$ be a fixed nonempty compact free $C_p$-space with a compatible
invariant metric.
Tensoring by $\tau_i=\Q[G_i/P_i]$, with its orthonormal form,
diagonal action and unchanged spatial labels, defines
\begin{equation}\label{alg:nu}
 \nu=p\,\mathrm{id}-\tau\otimes(-)
 \quad\text{on }L_*(\mathcal A_{\mathbf G}(T)).
\end{equation}

\begin{lemma}[Local triviality]\label{alg:localtrivial}
For compact $U\subset T/C_p$ contained in an open covering
trivialization, $\nu=0$ on $L_*(\mathcal A_{\mathbf G,T_U}(T))$,
where $T_U$ is its inverse image.
\end{lemma}
\begin{proof}
Choose a closed invariant neighborhood $L$ of $T_U$ inside that
trivialization, with sheet function $c:L\to C_p$ satisfying
$c(hx)=\bar h c(x)$, where $\bar h$ is the image of $h$ in $C_p$.
Labels of supported objects eventually lie in
$L$. The orthogonal equivariant map
\begin{equation}\label{alg:localisometry}
 e_a\otimes e_b\longmapsto e_b
 \quad\text{in copy }c(\ell_i(b))^{-1}a\in C_p
\end{equation}
has propagation zero. The sheet pieces of $L$ are compact and
positively separated, so controlled morphisms eventually preserve
them. Hence this is a natural isometry $\tau\otimes(-)\cong
\mathrm{id}^{\oplus p}$, inducing equality in ultimate $L$-theory.
\end{proof}

Choose a finite compact cover $U_1,\ldots,U_s$ of $T/C_p$ inside
trivializations, fixed independently of $i$, and put $T_a=T_{U_a}$.

\begin{lemma}[Nilpotence]\label{alg:nilpotence}
One has $\nu^s=0$ on $L_*(\mathcal A_{\mathbf G}(T))$.
\end{lemma}
\begin{proof}
Induct on $s$ over all such compact free spaces. Set
$A=T_1\cup\cdots\cup T_{s-1}$ and $B=T_s$.
The natural support Mayer--Vietoris sequence follows from
\cref{alg:filtration}. Since $\nu=0$ on $A\cap B$, $\partial(\nu x)=0$.
Thus $\nu x$ comes from $L_*(\mathcal A_A)\oplus L_*(\mathcal A_B)$.
Relabeling and induction kill the first summand by $\nu^{s-1}$;
local triviality kills the second by $\nu$. For $s>1$, applying
$\nu^{s-1}$ proves the claim; $s=1$ is local triviality.
\end{proof}

\begin{theorem}[Integer tail divisibility]\label{alg:divisibility}
Every $\alpha\in L_4(\mathcal A_{\mathbf G}(T))$ has signature tail
eventually divisible by $p$.
\end{theorem}
\begin{proof}
Extract forms $(V_i,b_i)$. Naturality and nilpotence imply
\begin{equation}\label{alg:factorannihilation}
 (p-\tau_i)^s[V_i,b_i]=0\quad\text{in }L_4^p(\Q[G_i])
 \quad\text{for all sufficiently large }i.
\end{equation}
For a generator $g_i$ of $G_i$, $\chi_{\tau_i}(g_i)=0$. Taking
signature characters gives $p^s\Sign_{G_i}(V_i,b_i)(g_i)=0$.
Apply \cref{alg:cyclotomic}.
\end{proof}

\section{Rational Poincar\'e corners}\label{alg:corners}

\subsection{Signs and rational structures}\label{alg:qmodel}
Complexes are homologically graded. We use
\begin{equation}\label{alg:dualconvention}
 (C^{N-*})_r=C_{N-r}^*,\qquad d_{C^{N-*},r}=(-1)^rd_C^*,\qquad
 (T\phi)_r=(-1)^{r(N-r)}\phi_{N-r}^*.
\end{equation}
The double shifted dual has differential $(-1)^{N+1}d_C$; its
identification with $C$ is $(-1)^{r(N-r)}=(-1)^{r(N+1)}$ in degree $r$.
Right slant sends $x_r\otimes y_{N-r}$ to
$\alpha\mapsto x_r\alpha(y_{N-r})$, without an extra sign.
For a bounded complex in a rational additive category with involution put
\[
 M(C)_k=\bigoplus_r\operatorname{Hom}(C_{k-r}^*,C_r),\qquad
 (\partial f)_r=d_Cf_{r+1}+(-1)^rf_rd_C^*,\qquad
 (Tf)_r=(-1)^{r(k-r)}f_{k-r}^*.
\]
Thus $T^2=1$, $T\partial=\partial T$, and conjugation by chain maps
commutes with $T$. For finite scalar complexes this is $C\otimes C$
under right slant. Exactness of rational invariants and coinvariants gives
\begin{equation}\label{alg:qgroups}
 Q^N(C)=H_N(M(C)^{C_2}),\qquad Q_N(C)=H_N(M(C)_{C_2}).
\end{equation}
The norm $1+T$ is an isomorphism with inverse $t\mapsto[t/2]$.
We write $\Sym=(1+T)/2$.
Consequently a symmetric structure has a strict symmetric representative
$\phi$, and its quadratic lift is $\psi_0=\phi/2$, $\psi_s=0$ for $s>0$.
If symmetric cycles differ by $\partial H$, replace $H$ by $(H+TH)/2$;
halving that primitive identifies their quadratic classes. The same
argument applies to the relative cone complexes, since their maps commute
with $T$ \cite[Proposition 3.3]{Ran80I}, \cite[Section 3]{Ran89}.
All these operations preserve propagation and exterior factorizations.

\subsection{Compatible idempotents and normalization}
Let a finite cyclic group $G$, of order $q$, act up to homotopy on
$(C,\phi)$, with
\begin{equation}\label{alg:grouprelations}
 A_gA_h\simeq A_{gh},\qquad A_1=1,\qquad
 A_g\phi\simeq\phi A_{g^{-1}}^*.
\end{equation}
Only an action in the chain-homotopy category is assumed; no higher
coherence data are used.
On $\mathcal C=C\otimes\Q[G]$, with strict left sheet action, set
\begin{equation}\label{alg:graphidempotent}
 \Phi=\phi\otimes1,\qquad
 E=q^{-1}\sum_gA_g\otimes R_{g^{-1}},
\end{equation}
where $R_{g^{-1}}$ sends sheet $b$ to $bg^{-1}$. Then
\begin{equation}\label{alg:graphrelations}
 E^2\simeq E,\qquad E\Phi\simeq\Phi E^*,\qquad
 E_{ab}=q^{-1}A_{a^{-1}b}.
\end{equation}
In the controlled application these equations, including their homotopies,
hold in the exterior quotient, as checked in \S\ref{sec:equivariant}.

\begin{lemma}[Compatible splitting]\label{alg:corner}
Let $\mathcal Q$ be rational additive with involution,
$\mathcal C\in[0,N]$, and
$\Phi:\mathcal C^{N-*}\to\mathcal C$ a strict symmetric Poincar\'e equivalence.
For an invertible self-adjoint central coefficient $q$ and a chain map
$E$ with homotopies $E^2\simeq E$, $E\Phi\simeq\Phi E^*$ in $\mathcal Q$,
there is a splitting in the bounded chain-homotopy category
$K^b(\Kar\mathcal Q)$,
\begin{equation}\label{alg:split}
 D\xrightarrow{\iota}\mathcal C\xrightarrow rD,\qquad
 \iota r\simeq E,\quad r\iota\simeq1,\qquad D\in[0,N].
\end{equation}
The strict symmetric Poincar\'e duality and a homotopy inverse are
\begin{equation}\label{alg:cornerform}
 \phi_D=q^{-1}r\Phi r^*,\qquad b_D=q\,\iota^*\Phi^{-1}\iota.
\end{equation}
Here $\Phi^{-1}$ denotes a chosen chain-homotopy inverse.
The resulting quadratic Poincar\'e complex determines a natural class
in $L_N(\mathcal Q)$.
\end{lemma}
\begin{proof}
With its split exact structure, $\Kar\mathcal Q$ has bounded derived
category $K^b(\Kar\mathcal Q)$, which is idempotent complete
\cite[Theorem 2.8]{BS01}. Hence $[E]$ splits. Outside $[0,N]$, the
identity homotopy $r\iota\simeq1$ makes the extremal differential split:
surjective at the bottom and injective at the top. Removing the resulting
contractible two-term summands puts $D$ in $[0,N]$.
Compatibility gives
\[
 \phi_Db_D\simeq r\Phi E^*\Phi^{-1}\iota\simeq rE\iota\simeq1,
 \qquad b_D\phi_D\simeq\iota^*\Phi^{-1}E\Phi r^*
                 \simeq\iota^*E^*r^*\simeq1.
\]
Conjugation preserves strict symmetry, so $\phi_D/2$ is its complete
quadratic structure. Two splittings are related by $r_2\iota_1$, an
equivalence carrying their symmetric structures to one another up to
homotopy. The primitive argument above and the Poincar\'e
mapping-cylinder cobordism \cite[Definition 3.10]{Ran89} identify their
ordinary $L$-classes in $\Kar\mathcal Q$. Finally pass to ultimate
decoration and use the natural isomorphism
$L_N(\mathcal Q)\cong L_N(\Kar\mathcal Q)$
\cite[Definition 17.7]{Ran92}. This also proves functoriality.
\end{proof}

For an asymptotic exterior category the splitting takes place once in
the sequence category. Its finite collection of maps and homotopies
therefore has asymptotically small representatives, with equations
factoring through exterior-supported objects.

\begin{lemma}[Scalar normalization]\label{alg:normalization}
For \eqref{alg:graphidempotent}, suppose a duality-preserving forgetting
functor retains a coordinate in which all $A_g$ and the homotopies in
\eqref{alg:grouprelations} are uniformly small. Its image of
\eqref{alg:cornerform} is homotopy-isometric to $(C,\phi)$.
\end{lemma}
\begin{proof}
Only after forgetting, let $j_1,\pi_1$ include and project the identity
sheet, and put $v=Ej_1$, $u=q\pi_1E$. Thus
$v_a=q^{-1}A_{a^{-1}}$, $u_b=A_b$, and
\begin{equation}\label{alg:uvrelations}
 uv\simeq1,\quad vu\simeq E,\quad uE\simeq u,\quad Ev\simeq v,
 \qquad u\Phi u^*=\sum_bA_b\phi A_b^*\simeq q\phi.
\end{equation}
Hence $w=u\iota$ and $z=rv$ are inverse equivalences and
\begin{equation}\label{alg:scalarisometry}
 w\phi_Dw^*\simeq q^{-1}uE\Phi E^*u^*
              \simeq q^{-1}u\Phi u^*\simeq\phi.
\end{equation}
The rational primitive argument identifies the quadratic classes;
the mapping-cylinder cobordism and ultimate completion invariance then
identify their $L$-classes. If $q=q_i$ varies, multiplication by
$q_i^{\pm1}$ has propagation zero, and the uniform entry and homotopy
bounds make $u,v,w,z$ valid sequence morphisms.
\end{proof}

\section{Localization boundaries}\label{sec:localization}
For a closed invariant cover $X=A\cup B$, put $Y=A\cap B$.
Support filtration, excision, and \cite[Theorem 4.2]{CP95} give
\begin{equation}\label{eq:ordinary-boundary}
 \partial_{A,B}:L_k(\mathcal A(X))\longrightarrow
                    L_{k-1}(\mathcal A_Y(X)).
\end{equation}
Here $\mathcal A=\mathcal A_{\mathbf G}$; the group sequence is unchanged.

\begin{proposition}\label{prop:exterior-boundary}
If $Z\subset A$ is closed invariant and $d(Z,B)>0$, there is a boundary
\begin{equation}\label{eq:exterior-boundary}
 \partial_{A,B}:L_k(\mathcal A(X)/\mathcal A_Z(X))
               \longrightarrow L_{k-1}(\mathcal A_Y(X)),
\end{equation}
natural for duality-preserving functors preserving these support categories.
\end{proposition}
\begin{proof}
Put $\mathcal Q=\mathcal A(X)/\mathcal A_Z(X)$ and
$\mathcal U=\mathcal A_A(X)/\mathcal A_Z(X)$.
Separation embeds $\mathcal A_B$ and $\mathcal A_Y$ fully faithfully in
$\mathcal Q$. Choose controlled representatives for finite families of
quotient maps and perform the based support decompositions before
quotienting; this verifies the support filtrations there. The same based
split as in \cref{alg:excision} gives the duality-preserving equivalence
\begin{equation}\label{eq:exterior-excision}
 \mathcal U/\mathcal A_Y\simeq\mathcal Q/\mathcal A_B.
\end{equation}
Indeed, on $A$-supported objects $I_B$ restricts to $I_Y$, whereas $I_Z$
already factors through $A$-supported objects; thus the respective ideals
are $I_Z+I_Y$ and $I_Z+I_B$. The based inclusion and projection are
adjoints and inverse in the quotient. The boundary is induced by
\[
 \Lspec(\mathcal Q)\longrightarrow\Lspec(\mathcal Q/\mathcal A_B)
 \simeq\Lspec(\mathcal U/\mathcal A_Y)
 \longrightarrow\Sigma\Lspec(\mathcal A_Y).
\]
Localization for the two filtrations gives Mayer--Vietoris and naturality.
The ordinary case is $Z=\varnothing$.
\end{proof}

\begin{lemma}\label{lem:union-boundary}
Let $V_A=(P\to D_A)$ and $V_B=(P\to D_B)$ be $k$-dimensional quadratic
Poincar\'e pairs supported on $A,B$, with oppositely oriented common
boundary $P$ supported on $Y$. Orient $P$ as the boundary of $V_A$.
For $V=V_A\cup_PV_B$, in the ordinary or exterior quotient setting,
\begin{equation}\label{eq:union-boundary}
 \partial_{A,B}[V]=[P].
\end{equation}
\end{lemma}
\begin{proof}
Use $\mathcal Q,\mathcal U$ above, with $Z=\varnothing$ in the ordinary
case. The relative-pair representative $P\to D_A$ in the cofiber of
$\Lspec(\mathcal A_Y)\to\Lspec(\mathcal U)$ has boundary $[P]$.
The relative-to-quotient comparison \cite[proof of Theorem 4.1]{CP95},
followed by ultimate decoration, sends it to $D_A$ with its full induced
quadratic structure. Modulo $\mathcal A_B$, $D_B$, $P$, and the
suspended interface of the union vanish degreewise, along with every
interface structure correction. The surviving complex and structure are
exactly those on $D_A$. Excision thus identifies the two representatives,
whose connecting image is $[P]$. Completion and the passage to ultimate
decoration preserve this pair comparison.
\end{proof}

\section{Local chain carriers}\label{sec:carriers}

Write $\varepsilon$ for geometric augmentation of zero chains.
We use the acyclic-carrier method to prove the required shrinking-mesh
estimates directly.  Its application to dual cells is the elementary PL
form of combinatorially controlled Poincar\'e--Lefschetz duality; compare
\cite[Proposition 5.3, Theorem 6.7 and Corollary 6.8]{Ran99}.

\begin{lemma}[Finite acyclic carriers]\label{lem:finite-acyclic-carrier}
Let $E$ be a finite based chain complex in nonnegative degrees.  Assign
to each basis element $e$ an augmented acyclic subcomplex $\Gamma(e)$
of a rational chain complex $F$, with augmentations agreeing on inclusions.
Assume $\Gamma(e')\subset\Gamma(e)$ whenever $e'$ occurs in $de$.
Two carried chain maps $u,v:E\to F$ agreeing in augmentation on degree-zero
basis elements admit a carried homotopy
\[
u-v=dH+Hd,\qquad H(e)\in\Gamma(e).
\]
Carried, augmentation-compatible prescribed degree-zero images likewise extend
to a carried chain map from an augmented source.
\end{lemma}
\begin{proof}
Induct on degree.  The residual $u(e)-v(e)-H(de)$ lies in $\Gamma(e)$,
is a cycle, and in degree zero has augmentation zero.  Fill it in that
carrier and use the filling as $H(e)$.  The map-existence induction fills
the cycle given by the already defined image of $de$.
\end{proof}

\begin{corollary}\label{cor:carrier-control}
If labels in $\Gamma(e)$ are within $a$ of the label of $e$, the constructed
map or homotopy has propagation at most $a$.  Relabeling by a fixed
continuous map $c$ on a compact metric space replaces this bound by
$\omega_c(a)$, where $\omega_c$ is a modulus of continuity for $c$.
Signed cellular bases use geometric augmentation.
\end{corollary}
\begin{proof}
The construction introduces no basis elements outside the assigned carriers.
\end{proof}

\begin{lemma}[Filling on a fixed triangulation]
\label{lem:unchanged-target-filling}
Let $K$ triangulate a compact $m$-manifold $W$ with mesh at most $\eta$.
Let $C$ be convex in one coordinate chart.  A simplicial $r$-cycle $\gamma$
whose full simplices lie in $C$ bounds a chain $q$ on the same $K$, with
full-simplex support in the closed $\eta$-neighborhood of $C$; for $r=0$
require augmentation zero.  For $r=m$ this forces $\gamma=0$.
\end{lemma}
\begin{proof}
The zero chain needs no filling.  Otherwise let $A$ be the subcomplex
generated by its support and cone its inclusion to a point in $C$.
This gives $F:|\operatorname{cone}A|\to C$ restricting to the inclusion
on $A$.  Subdivide only this cone and choose a simplicial approximation
$t:L\to K$ satisfying
$F(\operatorname{st}^{\circ}_Lv)\subset\operatorname{st}^{\circ}_Kt(v)$
\cite[Theorem 2C.1 and its proof]{Hat02}; write $S$ for subdivision.
On the base, $tS$ is carried in each original simplex and fixes vertices,
so $tS|_{C_*(A)}=1$: induction reduces the difference on an $r$-simplex
to a top-dimensional cycle in that simplex, which is zero for $r>0$.
Thus $q=tS(c\gamma)$ satisfies $dq=\gamma$, where $c$ is the augmented
cone operator.  For $z$ in the interior of a simplex $\lambda$ of $L$,
the approximation property says that every vertex of $t\lambda$ belongs
to the smallest target simplex containing $F(z)$.  That simplex meets
$C$ and has diameter at most $\eta$, proving the support claim.
In top degree $q=0$ because $C_{m+1}(K)=0$.
\end{proof}

\begin{lemma}[Relative dual-block carrier maps]
\label{lem:relative-block-carriers}
Let $K$ be a finite combinatorial triangulation of a compact oriented
$m$-manifold $A$ with
$\partial A$ a subcomplex and mesh $\eta$.  For nonboundary simplices
choose dual-block orientations with
\[
\partial D\sigma=(-1)^{\dim\sigma+1}
\sum_{\substack{\tau\supset\sigma\\\dim\tau=\dim\sigma+1}}
[\tau:\sigma]D\tau,\qquad D\rho=\epsilon_\rho[b_\rho]
\]
on top simplices, where $\epsilon_\rho$ is the fundamental coefficient
and $b_\rho$ the barycenter.
Let $D_{\mathrm{rel}}$ be the cellular subcomplex generated by these
blocks and put $J_\sigma=\bigcup_{\tau\supseteq\sigma,\,
\tau\not\subset\partial A}D\tau$.  Every $J_\sigma$ is augmented acyclic.
There are maps $S_{\mathrm{rel}}:D_{\mathrm{rel}}\to C_*(\operatorname{sd}K)$
and $R_{\mathrm{rel}}$ in the reverse direction, with
$R_{\mathrm{rel}}S_{\mathrm{rel}}=1$ and
$S_{\mathrm{rel}}R_{\mathrm{rel}}\simeq1$, all having propagation
$O_m(\eta)$ when cells are labeled by points of their geometric supports.
\end{lemma}
\begin{proof}
Use the dual orientations induced by the manifold and simplex orientations;
the displayed incidence formula follows.  On an interior codimension-one
simplex, the augmentation of its dual boundary is
$(-1)^m\sum_{\rho\supset\sigma}\epsilon_\rho[\rho:\sigma]=0$,
by the relative fundamental-chain equation.  Thus geometric augmentation
is a chain map, including the signed zero-cell convention.
All cofaces of a nonboundary simplex are nonboundary, giving the subcomplex.
If $\sigma$ is interior, $J_\sigma=D\sigma$ is a ball.  If $\sigma$ is
boundary, its link is an $l$-ball, $l=m-\dim\sigma-1$, with boundary the
link in $\partial A$.  Nonboundary cofaces correspond to nonboundary
faces $\lambda$ of this link.  With join orientations the map
\[
\lambda^*\longmapsto(-1)^{r(l+1)}D(\sigma*\lambda)
\]
identifies its relative shifted cochain complex in degree $r$ with
$C_r(J_\sigma)$.  Hence $H_r(J_\sigma)=\Q$ for $r=0$ and zero otherwise;
augmentation is nonzero on a top-coface dual vertex and identifies $H_0$
with $\Q$.
For a barycentric flag use $J_{\sigma_0}$, its minimum-face carrier.
These shrink on deleting the minimum and are unchanged on deleting other
faces.  The carrier lemma constructs $R_{\mathrm{rel}}$; take
$S_{\mathrm{rel}}$ to be cellular subdivision.  Augmentation fixes the
composite on zero cells; on a positive-dimensional $D\sigma$ the composite
lies in that dual ball and equals the identity by induction and absence
of top-dimensional cycles there.  The other composite
and the identity are carried in $\operatorname{sd}(\operatorname{st}_K
\sigma_0)$, so the carrier lemma supplies their homotopy.
Closed stars are contractible and have diameter at most $2\eta$,
giving every asserted bound independently of valence and rank.
\end{proof}

\begin{lemma}[Controlled cap--block comparison]
\label{lem:cap-block-comparison}
In the setting of Lemma~\ref{lem:relative-block-carriers}, put
$R=C_*(K,\partial K;\Q)$ and write $\delta$ for its shifted-dual
differential $(-1)^rd^*$. Let $F:R^{m-*}\to C_*(K)$ be a chain map
carried, on $\sigma^*$, in the closed star $\operatorname{st}_K\sigma$.  Suppose
$\varepsilon F(\rho^*)=\epsilon_\rho$ on each top simplex.
Then $F$ has a chain homotopy inverse and both inverse homotopies of propagation
$O_m(\eta)$.  More precisely, it is carried-homotopic to $F_0$, with
homotopy inverse $B$, where
\[
F_0=R_KS_{\mathrm{rel}}\mathcal J,
\qquad B=\mathcal J^{-1}R_{\mathrm{rel}}S_K,
\qquad\mathcal J_r(\sigma^*)=(-1)^{r(m+1)}D\sigma,
\]
where $S_K$ is subdivision and $R_K$ a simplex-carried inverse.
Fixed continuous relabelings give bounds tending to zero with $\eta$.
\end{lemma}
\begin{proof}
The incidence formula makes $\mathcal J$ a chain isomorphism.
The carrier lemma on maximum-face carriers constructs $R_K$, with
$R_KS_K=1$ by top-chain uniqueness in a simplex; it also gives
$S_KR_K\simeq1$ in carriers $C_*(\operatorname{sd}\sigma_r)$ on flags
with maximum $\sigma_r$.  Together with the preceding lemma these
identities prove $F_0B\simeq1$ and $BF_0\simeq1$ by a fixed number of
compositions.  Both $F,F_0$ have the same top augmentations and closed-star
carriers.  The dual differential uses cofaces, and
$\operatorname{st}_K\tau\subset\operatorname{st}_K\sigma$ for
$\tau\supset\sigma$, so the carrier lemma gives $F-F_0=dH+H\delta$.
Add $HB$ and $BH$ to the inverse homotopies for $F_0$.
All carriers have diameter at most $2\eta$; fixed compositions give
$C_m\eta$ and relabeling gives $\omega_c(C_m\eta)$.
\end{proof}

\begin{corollary}[A common simplicial diagonal]
\label{cor:common-diagonal-control}
In the setting of Lemma~\ref{lem:relative-block-carriers}, let
$j$ be the boundary inclusion.
Use one ordered Alexander--Whitney diagonal $\Delta$ on $K$ and its subcomplexes.
Choose oriented fundamental chains satisfying $dz_A=jz_{\partial A}$, and put
$t_A=\frac12(1+T)\Delta z_A$ and
$t_{\partial A}=\frac12(1+T)\Delta z_{\partial A}$.
Then $dt_A=j^{\otimes2}t_{\partial A}$ exactly. The right slant of $t_A$
on relative cochains is a chain equivalence with inverse data of propagation
$O_m(\eta)$, also when the boundary is empty.
\end{corollary}
\begin{proof}
The diagonal restricts to subcomplexes.  Its boundary equation gives
\[
d\phi_{A,r}-(-1)^r\phi_{A,r-1}d^*
=j\phi_{\partial A,r-1}j^*.
\]
The right side vanishes on relative cochains.  Each diagonal term consists of faces of one simplex, so slant
on $\sigma^*$ has closed coface-star support.  On a top simplex it is
$\epsilon_\rho([v_0]+[v_m])/2$, of augmentation $\epsilon_\rho$.
Apply the cap--block comparison.
\end{proof}

\section{Geometric data and quotient representatives}\label{geo:fixed-data}

Suppose that $\Zp$ acts effectively on a connected boundaryless
$n$-manifold, $n\ge1$, and put $H_j=p^j\Zp$.
First note that, for any compact group $K$ acting on a connected manifold,
every countable union of compact sets lies in a $K$-invariant connected
second-countable open submanifold. Indeed, every compact set lies in a
relatively compact connected open set covered by finitely many charts:
connect a finite chart cover by finitely many paths. Choose such sets
$V_j$ inductively with $K\overline{V_j}\subset V_{j+1}$, including one
prescribed compact set at each stage; their union has the required
properties. For $K=\Zp$, include a point moved by $p^j$ for every $j\ge0$.
The restriction remains effective, so we henceforth work on this
second-countable open submanifold.

\begin{lemma}\label{geo:free-point}
There is a point with trivial stabilizer.
\end{lemma}
\begin{proof}
Every nonzero closed subgroup of $\Zp$ is an $H_j$: an element of
minimum valuation $j$ has integer multiples dense in $H_j$. If all stabilizers
were nonzero, the discrete subgroup $\Z$ would have finite point orbits.
By \cite[Theorem 2]{Rob75} its image would be finite. Some $m>0$ would
act trivially; the closed kernel would then contain
$\overline{m\Z}=H_{v_p(m)}$, contradicting effectivity.
\end{proof}

Center a chart at this point, with image containing $\overline{B(0,4)}$.
Joint continuity gives $H=H_j$ and $0<D<1/64$ such that
\begin{equation}\label{geo:small-H}
 |hx-x|\le D\qquad(h\in H,\ x\in\overline{B(0,2)}).
\end{equation}
The invariant saturation $\Omega=HB(0,2)$ lies in $B(0,2+D)$ and has
$H$-displacements at most $2D$. Its coordinate average with respect to
normalized Haar measure $\mu$
\begin{equation}\label{geo:average}
 F_0(x)=\int_H hx\,d\mu(h)
\end{equation}
is continuous, invariant, and satisfies
\begin{equation}\label{geo:average-bounds}
 |F_0(x)-x|\le2D,\qquad F_0(x)=x\text{ on }\Fix(H).
\end{equation}
Choose $0<\rho<\rho_+<1/4$ with
$\overline{B(0,\rho_+)}\cap\Fix(H)=\varnothing$. Then
\begin{equation}\label{geo:Zplus}
 Z^+=F_0^{-1}(\overline{B(0,\rho_+)})
\end{equation}
is compact, invariant, and disjoint from $\Fix(H)$: it is a closed
preimage contained in $B(0,\rho_++2D)\Subset\Omega$.
Choose $s_\rho:\R^n\to S^n$ restricting to an orientation-preserving
homeomorphism from $B(0,\rho)$ to $S^n\setminus\{\infty\}$, and constant
$\infty$ elsewhere. Define
\begin{equation}\label{geo:f}
 f=s_\rho F_0\text{ on }\Omega,\qquad f=\infty\text{ elsewhere}.
\end{equation}
Its nonconstant locus lies in $B(0,\rho+2D)$, so it extends to the
compactified coordinate sphere. The homotopy
\begin{equation}\label{geo:pinch-homotopy}
 s_\rho((1-v)F_0(x)+vx),\quad 0\le v\le1,
\end{equation}
extended by $\infty$, has the same compact support bound. Hence $f$ has
degree one and $Z^+$ is nonempty.

Put $H'=H_{j+1}$. Every proper closed subgroup of $H$ lies in $H'$, so
$C_p=H/H'$ acts freely on the compact metric space $Z^+/H'$. Its orbit
map is a covering: a neighborhood can be shrunk to have disjoint
translates. A finite trivializing cover and a subordinate partition
of unity give, by multiplying sheet coordinates by the square roots
of the partition functions, an equivariant detector
\begin{equation}\label{geo:detector}
 \widetilde\rho:Z^+\longrightarrow T=S^{2m-1}\subset\C^m,
 \qquad\widetilde\rho(hx)=\chi(h)\widetilde\rho(x),
\end{equation}
where $\chi:H\to H/H'=C_p\subset S^1$ is the quotient character.
Scalar multiplication makes $T$ a
fixed compact free isometric $C_p$-space.
Choose a round target metric with $y=s_\rho(0)$ antipodal to $\infty$,
and fix
\begin{equation}\label{geo:buffers}
 0<r<R<t<u<d(y,\infty),\quad
 Z_0=\{z:d(y,z)\ge R\},\quad Z=T\times Z_0.
\end{equation}
For every $v<d(y,\infty)$, the compact invariant set
$\mathcal K_v=f^{-1}(\overline{B(y,v)})$ lies in $\operatorname{int}Z^+$.
Only now choose
\begin{equation}\label{geo:finite-quotients}
 K_i=H_{\ell_i}\subset H',\quad\ell_i\to\infty,\qquad
 G_i=H/K_i,\quad P_i=H'/K_i.
\end{equation}
Write $\chi$ also for the induced character on each $G_i$.
The $K_i$-displacements on $H\overline{B(0,2)}$ tend uniformly to zero;
$T$, the detector, and all buffers remain fixed.

\section{Common chain models}\label{geo:chain-models}

Put $W=S^n\times\CP^2$, $N=n+4$. The sphere factor is the compactified
coordinate space. Fix a product metric, a finite smooth atlas with
uniform chart buffers and Lipschitz constants, and a common continuity
modulus $\omega$ for the compact map families below. Choose a
combinatorial triangulation compatible with the later finite PL cutting flag and
barycentric refinements $Y_i$ of mesh $\eta_i\to0$. Write
$D_i=C_*(Y_i;\Q)$ and let $z_i$ be its oriented fundamental cycle;
bases are simplices labeled at their barycenters. All control maps
ignore $\CP^2$, and support estimates concern whole closed simplices.

\begin{lemma}\label{geo:closed-models}
The $D_i$ have strict symmetric cap equivalences
$\phi_i:D_i^{N-*}\to D_i$, inverses $b_i:D_i\to D_i^{N-*}$,
and both inverse homotopies of propagation $O_N(\eta_i)$. An
equicontinuous family of self-maps has chain approximations on these
unchanged complexes, with uniform full-simplex graph error
\begin{equation}\label{geo:mu}
\mu_i=O(\eta_i+\omega(C\eta_i))\longrightarrow0.
\end{equation}
\end{lemma}
\begin{proof}
The cap assertion is \cref{cor:common-diagonal-control} with empty
boundary. Subdivide the source enough for the target open-star cover
and the common modulus, and take $a_h=t_hS$ by simplicial approximation
\cite[Theorem 2C.1]{Hat02}. The identity approximation is exactly $1$
using the last-vertex map. Each image simplex lies within one target
mesh of $h$ of a point of its original source simplex; equicontinuity
gives \eqref{geo:mu}.
\end{proof}

Choose a continuous cutoff $0\le\theta\le1$, equal to $1$ on $B(0,1)$
and $0$ off $B(0,3/2)$, and complete each $h\in H$ to a sphere map
\begin{equation}\label{geo:completed-map}
 \widehat h(x)=x+\theta(x)(hx-x),
\end{equation}
using the identity elsewhere. These maps are equicontinuous and degree
one, by scaling $\theta$, and
\begin{equation}\label{geo:f-invariance}
 f\widehat h=f:
\end{equation}
inside $B(0,1)$ use invariance; outside it both points miss the
nonconstant locus. Choose representatives $h_g$ of $G_i$, $h_1=1$, and
write $h_gh_h=h_{gh}k_{g,h}$ with $k_{g,h}\in K_i$. The tracks
\begin{equation}\label{geo:raw-product-track}
 L_{g,h}(x,v)=(1-v)\widehat h_g\widehat h_h(x)+v\widehat h_{gh}(x)
\end{equation}
are constant at infinity. For $|x|\le1/2$ they lie within
$\beta_i\to0$ of $h_{gh}x$, uniformly; for $|x|\ge1/2$ they lie within
$2D$ of $x$ and have $f$-value $\infty$.

A sequence of matrices has graph type $g$ if its global $f$-propagation
tends to zero and, for every fixed $v<d(y,\infty)$, its entries $x\to x'$ with
$x\in\mathcal K_v$ satisfy
$d_T(\widetilde\rho(x'),\chi(g)\widetilde\rho(x))\to0$ uniformly.
Their target labels are defined eventually. A larger fixed $v'>v$ gives
the same estimates for transposes and products of bounded length: the transpose
has type $g^{-1}$, and graph types multiply.
Apply \cref{geo:closed-models} to $\widehat h_g\times1$ to obtain
$a_{g,i}$, $a_{1,i}=1$.
Equation \eqref{geo:f-invariance}, detector equivariance, and the graph
error give type $g$. Degree one and the absence of $(N+1)$-chains give
\begin{equation}\label{geo:fundamental-cycle}
 a_{g,i}z_i=z_i.
\end{equation}
We now suppress the index $i$.

\begin{lemma}\label{geo:product-homotopies}
On the same complexes there are homotopies
\begin{equation}\label{geo:product-equation}
 a_ga_h-a_{gh}=dB_{g,h}+B_{g,h}d
\end{equation}
of graph type $gh$, uniformly in $i,g,h$.
\end{lemma}
\begin{proof}
Induct over simplex dimension, filling the residual cycle
$\gamma_\sigma=a_ga_h\sigma-a_{gh}\sigma-B(\partial\sigma)$ using
\cref{lem:unchanged-target-filling}; in degree zero it has augmentation
zero. Take first-coordinate convex hulls of
$L_{g,h}(|\sigma|\times[0,1])$, and small coordinate-convex carriers for the
second projection. In a product chart, take convex hulls including full
endpoint chains and already fixed face fillings. With
$\delta_i=\eta_i+\mu_i+\omega(C\eta_i)+\omega(C\mu_i)+\beta_i\to0$,
the full supports stay within widths $w_r$ of the nested track hulls
and second-coordinate projections, where
$w_{-1}=0$, $w_r\le C(w_{r-1}+\delta_i)$: coordinate changes cost a
fixed factor and filling adds one target mesh. Thus $w_r\to0$ uniformly
for $r\le N$, independently of all ranks and auxiliary subdivisions.
These convex carriers lie in one product chart; near infinity use
small charts in the identity region. The top residual vanishes by the
same filling lemma, so no extra top chain is required.

For a source barycenter with $|x|\le3/8$, every filling lies within
$o(1)$ of $h_{gh}x$. Equivariance of the detector on each $\mathcal K_v$ gives
type $gh$. For $|x|\ge3/8$ its support is within $2D+o(1)$ of $x$;
the fixed margin $3/8-\rho-4D>1/16$ makes its $f$-value $\infty$.
Uniform continuity finishes both control estimates.
\end{proof}

The same filling recursion, now on the based source $D^{N-*}$, gives
a local primitive $U_g$ with
\begin{equation}\label{geo:isometry-homotopy}
 a_g\phi a_g^*-\phi=dU_g+U_gd_{D^{N-*}}.
\end{equation}
Indeed, an entry of $a_g^*$ selects a source cell mapping near its
given target cell; $\phi$ moves within $O(\eta_i)$, and $a_g$ maps back
near that target cell. Thus both endpoint maps are local with error
$O(\mu_i+\omega(C\eta_i)+\eta_i)\to0$. Their degree-zero
augmentations agree, since simplicial approximations preserve
geometric augmentation ($\varepsilon a_g=\varepsilon$): for
$\alpha\in D_N^*$,
$\varepsilon\phi(\alpha)=\alpha(z_i)$, while
$\varepsilon a_g\phi a_g^*(\alpha)=\alpha(a_gz_i)=\alpha(z_i)$.
For $\sigma^*$ take a chart-convex hull of its endpoint supports and
the carriers and fillings already assigned to cofaces occurring in
the shifted dual differential, enlarging by one mesh after filling.
These geometric carriers are nested under cofaces.
Every coface meets $\sigma$, so the same width induction keeps them
within $o(1)$ of $\sigma$. The top residual again vanishes.
Put $a=a_g$, $b'=a_{g^{-1}}$. For $b'a-1=dT+Td$ set
$K_r=(-1)^{r+1}T_{N-r-1}^*$. The identity
$a\phi-\phi b'^*=a\phi(1-a^*b'^*)+(a\phi a^*-\phi)b'^*$ yields
\begin{equation}\label{geo:adjoint-homotopy}
 V_g=-a\phi K+U_gb'^*,\qquad
 dV_g+V_gd_{D^{N-*}}=a_g\phi-\phi a_{g^{-1}}^*.
\end{equation}
Since $T,K,U_g$ have type $1$ and $b'^*$ type $g$, $V_g$ has type $g$.
Fix one $\epsilon_i\to0$ bounding all these $f$ and detector errors,
cell oscillations, and local duality propagations on the fixed buffers.

\section{Compression and exterior control}\label{geo:compression}

Keep the fixed detector neighborhood and $R<t<u$. All primitive local
and graph errors below are bounded uniformly by $\epsilon_i\to0$;
cell oscillations of $f$ are bounded by $\kappa_i\to0$.
Let $F_i\subset D_i$ be spanned by simplices whose closed simplices miss
$f^{-1}(\overline{B(y,t)})$. This is a subcomplex, since faces of a
missing simplex also miss that set. Put
\begin{equation}\label{geo:compression-data}
 C=D/F,\qquad \pi:D\to C,\qquad \sigma:C\to D,
 \qquad Q=1-\sigma\pi,\qquad \widehat d=\pi d\sigma,
\end{equation}
where $\sigma$ is the degreewise based section. Thus $\pi dQ=0$.
Call the simplices surviving in $C$ retained and those in $F$ far.
Retained labels have sphere radius $d(y,f(x))$ at most $t+\kappa_i$;
far labels have radius greater than $t$. Only retained matrices and the restricted
witnesses below need detector labels; the following computations on
$D_i$ are raw matrix computations.

\begin{lemma}[Exterior band]\label{geo:band-factorization}
Consider words in the preceding primitive maps and homotopies, their
shifted duals, and the based projections and sections. A
retained-to-retained word of fixed length at most $L$ containing
$Q$ or $Q^*$ factors through a based module labeled in
\begin{equation}\label{geo:band}
 t-L\epsilon_i\le d(y,f(x))\le t+\kappa_i+L\epsilon_i.
\end{equation}
Local words have controlled scalar factors; graph words have controlled
equivariant factors after compensating sheet translations. Their cut
objects are eventually supported in the fixed exterior $Z$.
\end{lemma}
\begin{proof}
At a far projection restrict to basis elements reached from retained
sources and contributing to retained targets, and cut the matrix product
there. This is an exact factorization. The far label has radius greater
than $t$, and a retained endpoint and at most $L$ sphere errors give the
upper bound in \eqref{geo:band}. Every contributing intermediate label
also has radius at most this upper bound, eventually less than $u$;
all required detector labels therefore exist in the fixed buffered domain.
The cut lies outside $R$. Local propagation is at most $L\epsilon_i$.
For graph words use all sheets, label $(x,s)$ by
$(\chi(s)\widetilde\rho(x),f(x))$, and accompany a graph-$g$ primitive by
$s\mapsto sg^{-1}$. Its detector displacement cancels because
$\chi(sg^{-1})\chi(g)=\chi(s)$. Apply this separately to the prefix
and suffix. Their intermediate basis is a free $G_i$-set and both factors
are controlled. Sums factor through direct sums of the cut objects;
their propagation is bounded by the maximum summand bound.
\end{proof}

Let $b_i:D_i\to D_i^{N-*}$ be the local cap inverse supplied by
\cref{geo:closed-models}, and put
\begin{equation}\label{geo:compressed-maps}
 A_g=\pi a_g\sigma,\qquad
 \widehat\phi=\pi\phi\pi^*,\qquad
 \widehat b=\sigma^*b\sigma.
\end{equation}
For a continuous $c:W\to S^n$, denote by
$\mathfrak m_c\in L_N(\mathcal B_{\mathbf1,Z_0}(S^n))$ the scalar
exterior class of $(D_i,\phi_i/2)$ with labels $c$.

\begin{proposition}\label{geo:compressed-realization}
The compressed complex is Poincar\'e in
$\mathcal B_{\mathbf1,Z}(T\times S^n)$, with quadratic structure
$\widehat\phi/2$. Its group and adjoint homotopies have the stated
uniform graph control modulo controlled exterior witnesses, and $A_1=1$.
After forgetting $T$, its class is $\mathfrak m_f$.
\end{proposition}
\begin{proof}
For a degree-compatible raw operator put $\widehat\alpha=\pi\alpha\sigma$.
Insertion of $\sigma\pi=1-Q$ gives
\begin{align}
 \widehat\alpha\widehat\beta-\widehat{\alpha\beta}
   &=-\pi\alpha Q\beta\sigma,\label{geo:compression-product}\\
 \widehat dA_g-A_g\widehat d
   &=\pi a_gQd\sigma,\label{geo:compression-chain}\\
 A_gA_h-A_{gh}-\widehat d\widehat B-\widehat B\widehat d
   &=-\pi a_gQa_h\sigma+\pi B_{g,h}Qd\sigma,
   \quad\widehat B=\pi B_{g,h}\sigma.
   \label{geo:compression-relation}
\end{align}
The band lemma kills all defects after compensation. The map
$\widehat\phi$ is a strict symmetric chain map because $\pi$ is a chain
map. For its inverse, let $\delta,\widehat\delta$ be the shifted-dual
differentials of $D,C$, and put
$e=d\sigma-\sigma\widehat d=Qd\sigma$ and
$e^\vee=\widehat\delta\sigma^*-\sigma^*\delta$, so
$(e^\vee)_r=(-1)^{r+1}e_{N-r+1}^*$ factors through $Q^*$.
Then
\[
 \widehat\delta\widehat b-\widehat b\widehat d
 =e^\vee b\sigma+\sigma^*be.
\]
Compress the local inverse homotopies
$\phi b-1=dH+Hd$ and $b\phi-1=\delta H'+H'\delta$:
\begin{align*}
 \widehat\phi\widehat b-1
 &=\widehat d(\pi H\sigma)+(\pi H\sigma)\widehat d
      +\pi He-\pi\phi Q^*b\sigma,\\
 \widehat b\widehat\phi-1
 &=\widehat\delta(\sigma^*H'\pi^*)
      +(\sigma^*H'\pi^*)\widehat\delta
      -e^\vee H'\pi^*-\sigma^*bQ\phi\pi^*.
\end{align*}
All residuals are local bounded-length words containing $Q$ or $Q^*$.
Similarly, for $\widehat V_g=\pi V_g\pi^*$,
\[
 A_g\widehat\phi-\widehat\phi A_{g^{-1}}^*
 -\widehat d\widehat V_g-\widehat V_g\widehat\delta
 =-\pi a_gQ\phi\pi^*
       +\pi\phi Q^*a_{g^{-1}}^*\pi^*.
\]
These residuals have graph type $g$ and vanish after compensation by
the same lemma. Finally $A_1=\pi\sigma=1$ and
\begin{equation}\label{geo:compressed-orientation}
 A_g\pi z_i-\pi z_i=-\pi a_gQz_i,
\end{equation}
whose output misses $\overline{B(y,R)}$ eventually.
After forgetting $T$, the entire $D_i$ is labeled by $f$ and $F_i$ is
$Z_0$-supported. Consequently $Q$ and $e$ vanish in the scalar exterior
quotient, so $\pi,\sigma$ are inverse chain equivalences. Transport of
$\phi/2$ along $\pi$ is exactly $\widehat\phi/2$.
\end{proof}

\section{The graph corner and its scalar image}\label{sec:equivariant}

Give $\mathcal C_i=C_i\otimes\Q[G_i]$ the free left sheet action and
labels as follows. Let $x_c$ be the first coordinate of the barycenter
of a retained simplex $c$. Then
\begin{equation}\label{eq:sheet-labels}
 \ell(c\otimes s)=
 (\chi(s)\widetilde\rho(x_c),f(x_c)).
\end{equation}
Copying local scalar matrices sheetwise preserves duality, exterior
support, and propagation. Set
\begin{equation}\label{eq:graph-idempotent}
 \Phi=\widehat\phi\otimes1,\qquad
 E_i=q_i^{-1}\sum_{g\in G_i}A_{g,i}\otimes R_{g^{-1}},
 \qquad q_i=|G_i|,\quad q=(q_i).
\end{equation}
The detector cancellation in the band lemma applies to the maps and
all product and adjoint homotopies. In particular
$q_i^{-2}\sum_{g,h}\widehat B_{g,h}\otimes R_{(gh)^{-1}}$
is a controlled product homotopy modulo the direct sum of the exterior
witnesses. Uniform control, rather than a bound on the number of
summands, gives
\begin{equation}\label{eq:controlled-idempotent}
 E^2\simeq E,\qquad E\Phi\simeq\Phi E^*
 \quad\text{in }K^b(\mathcal B_{\mathbf G,Z}(T\times S^n)).
\end{equation}
By \cref{alg:corner}, split $E$ in the Karoubi-completed homotopy category:
\begin{equation}\label{eq:normalized-corner}
 D\xrightarrow{\iota}\mathcal C\xrightarrow{r}D,
 \qquad \phi_D=q^{-1}r\Phi r^*.
\end{equation}
Its quadratic structure is $\phi_D/2$. Ultimate completion invariance
gives
\begin{equation}\label{eq:d-class}
 \mathfrak d\in L_{n+4}(\mathcal B_{\mathbf G,Z}(T\times S^n)).
\end{equation}

Let $U$ forget equivariance and $T$, retaining sphere control.
Every $A_g$ and every relation homotopy is uniformly small in this
coordinate. Scalar normalization (\cref{alg:normalization}) and
the last assertion of \cref{geo:compressed-realization} give
\begin{equation}\label{eq:scalar-image}
 U_*\mathfrak d=\mathfrak m_f
 \quad\text{in }L_{n+4}(\mathcal B_{\mathbf1,Z_0}(S^n)).
\end{equation}

\section{Evaluation of the iterated boundary}\label{bd:scalar}
Put $N=n+4$ and $W=S^n\times\CP^2$, with product orientation;
all sphere control ignores $\CP^2$. Fix
\[
 A=T\times(S^n\setminus B(y,r)),\quad
 B=T\times\overline{B(y,r)},\quad Y=T\times S_r(y).
\]
Here $S_r(y)=\partial B(y,r)$. The exterior $Z\subset A$ is positively separated from $B$.
Apply \cref{prop:exterior-boundary}, relabel its target onto $Y$, and
successively cut the sphere interfaces into two fixed hemispheres.
After $n$ boundaries the support is $T\times S^0$; project onto one
component. This gives
\begin{equation}\label{bd:Delta}
 \Delta_T:L_N(\mathcal B_{\mathbf G,Z}(T\times S^n))
       \longrightarrow L_4(\mathcal A_{\mathbf G}(T)).
\end{equation}
Every cut, relabeling, orientation, and component is fixed before the
index varies. The same choices give $\Delta_{\mathbf1}$ without $T$
and equivariance, with target $L_4(\mathcal V_{\mathbf1})$.
Forgetting equivariance and $T$, while retaining sphere control, gives
\begin{equation}\label{bd:naturality}
 U_*\Delta_T=\Delta_{\mathbf1}U_*.
\end{equation}
These boundaries also commute with $\tau\otimes(-)$, by localization
naturality.

\begin{lemma}[Controlled relative duality]\label{bd:relative}
Let $A$ be a compact oriented combinatorial $m$-manifold with boundary.
Use a finite combinatorial triangulation, with boundary a subcomplex,
and one ordered
diagonal. Write $C=C_*(\partial A;\Q)$, $D=C_*(A;\Q)$,
$j:C\to D$, and $J=\operatorname{Cone}(j)$. Fundamental chains
$dz_A=jz_{\partial A}$ give invariant tensors
\[
 t_D=\operatorname{Sym}\Delta z_A,\quad
 t_C=\operatorname{Sym}\Delta z_{\partial A},\quad
 dt_D=j^{\otimes2}t_C,\quad dt_C=0.
\]
Write $\phi_D,\phi_C$ for their right slants, which give the relative duality map
\begin{equation}\label{bd:Psi}
 \Psi_r(\alpha,\beta)=\phi_{D,r}\alpha
       +(-1)^{r+1}j\phi_{C,r}\beta:
       J^{m-*}_r\longrightarrow D_r.
\end{equation}
For mesh $\eta$, this map, an inverse, and both inverse homotopies
have geometric propagation $O_m(\eta)$. Any fixed continuous labeling
therefore gives an asymptotically controlled Poincar\'e pair, whose
full quadratic structure is $(t_D,t_C)/2$.
\end{lemma}
\begin{proof}
The tensor equation and
\[
 d_{J^{m-*},r}=(-1)^r
 \begin{pmatrix}d_D^*&0\\j^*&-d_C^*\end{pmatrix}
\]
make \eqref{bd:Psi} a chain map: its two relative-boundary terms cancel.
Let $\mathcal R=D/C$ and $q:J\to\mathcal R$ be the quotient.
For the based section $s:\mathcal R\to D$, write
$je=d_Ds-s\bar d$. Then $\tau=(s,-e)$ is a chain section of $q$;
$\ker q=\operatorname{Cone}(1_C)$ has contraction $k(x,c)=(0,x)$.
Thus
\begin{equation}\label{bd:cone-contraction}
 dh+hd=1-\tau q,\qquad h=k(1-\tau q).
\end{equation}
All these maps use only a simplex and its faces. By
\cref{cor:common-diagonal-control}, $F=\Psi q^*$ has a local cap inverse
$b$ and local inverse homotopies. Consequently $q^*b$ is an inverse
of $\Psi$: use $\Psi\simeq F\tau^*$ and
$q^*\tau^*\simeq1$, obtained by shifted dualizing
\eqref{bd:cone-contraction}. Only finitely many local compositions are
used, so their propagation is $O_m(\eta)$, independently of rank and
valence. Apply the labeling's continuity modulus. Taking halves gives
the full relative quadratic structure by \S\ref{alg:qmodel}.
\end{proof}

\begin{lemma}[Controlled unions]\label{bd:union}
A cut of a closed oriented combinatorial manifold $W$ into two manifolds
along a common boundary, using one triangulation and diagonal, gives
controlled Poincar\'e pairs whose union has the closed manifold's full
rational quadratic structure.
\end{lemma}
\begin{proof}
Write the chains as $A,B,P$, with inclusions $j_A,j_B$, and choose
$z_W=z_A+z_B$, $dz_A=j_Az_P$, $dz_B=-j_Bz_P$.
On $V=\operatorname{Cone}((j_A,-j_B))$, let $e_A,e_B$ include the
two tops, put $a=e_Aj_A$, $b=e_Bj_B$, and let $s$ include the
suspended interface. Then $ds+sd=a-b$. With tensor Koszul signs, set
\[
 K=s\otimes a+b\otimes s,\qquad
 K_{\rm sym}=(K+T_VKT_P)/2,
\]
\[
 t_V=e_A^{\otimes2}t_A+e_B^{\otimes2}t_B-K_{\rm sym}t_P.
\]
The identity $dK+Kd=a^{\otimes2}-b^{\otimes2}$ proves $dt_V=0$.
This is the Poincar\'e union of \cite[\S3]{Ran80I}.
The collapse $p:V\to C_*(W;\Q)$ has kernel
$\operatorname{Cone}(1_P)$ and satisfies
\begin{equation}\label{bd:union-tensor}
 (p\otimes p)t_V=t_W.
\end{equation}
Use the kernel's standard contraction $k$.
For a based section $r$ of $p$, $\ell=r-k(dr-rd)$ is a chain
section and $k(1-\ell p)$ contracts $1-\ell p$.
These maps, the interface correction, and the cap inverse of $W$
are local. Hence $p$ is a controlled chain equivalence identifying
the full symmetric structures; their halves identify the quadratic
structures. The propagation bounds are $O_m(\eta)$ before labeling
and tend to zero after any fixed continuous labeling.
\end{proof}

\begin{lemma}[Scalar control homotopies]\label{bd:homotopy}
Let $Y_i$ be the fixed ordered combinatorial triangulations of a compact closed
oriented PL $m$-manifold $W$, of mesh $\eta_i\to0$, and let
$c:W\times[0,1]\to X$ be a fixed continuous map to a compact metric
space. The complexes $C_*(Y_i;\Q)$ with their prescribed fundamental
structures and endpoint labels $c_0,c_1$ are rational quadratic
Poincar\'e cobordant. Their $L$-classes remain equal in every fixed
exterior quotient.
\end{lemma}
\begin{proof}
Triangulate the product prisms with time mesh $\lambda_i\to0$, keeping
both endpoint triangulations, vertex orders, and fundamental structures
exactly those of $Y_i$. Choose one vertex order restricting to the
prescribed endpoint orders and use its ordered diagonal. The cylinder
fundamental chain
\[
 Z_i=(-1)^m z_i\times[0,1],\qquad
 dZ_i=j_1z_i-j_0z_i
\]
gives the exact relative tensor equation with boundary $t_1-t_0$.
Label simplex barycenters by $c$. By \cref{bd:relative}, the pair's
duality, inverse, and inverse homotopies have propagation at most
$\omega_c(C_m(\eta_i+\lambda_i))\to0$ and degree range $[0,m+1]$.
Taking halves gives the endpoint quadratic structures. Apply the
exterior quotient functor to this completed cylinder; no condition on
the homotopy's behavior on the exterior is required.
See \cite[Propositions 6.1--6.3]{Ran80II} for the relative symmetric
construction.
\end{proof}

\begin{proposition}[Boundary evaluation]\label{bd:pinch-value}
For one constant $\varepsilon\in\{1,-1\}$ determined by the fixed
orientation and component choices,
\begin{equation}\label{bd:CP2}
 \Delta_{\mathbf1}\mathfrak m_f
   =\varepsilon[C_*(\CP^2;\Q),\phi_{\CP^2}/2],\qquad
 \sigma_{\rm tail}(\Delta_{\mathbf1}\mathfrak m_f)
   =[(\varepsilon,\varepsilon,\ldots)].
\end{equation}
\end{proposition}
\begin{proof}
The fixed scalar homotopy from $f$ to $s_\rho$ and
\cref{bd:homotopy} give
\begin{equation}\label{bd:pinch-homotopy}
 \mathfrak m_f=\mathfrak m_{s_\rho}.
\end{equation}
The degree-one sphere map $s_\rho$ is homotopic to an
orientation-preserving homeomorphism $g$ carrying the fixed source PL
disk $D$ and hemisphere flag, chosen before the $Y_i$, to the target ball
and flag; degree
classifies sphere maps for $n\ge1$ \cite[Corollary 4.25]{Hat02}.
Keep every target cut fixed and apply \cref{bd:homotopy} again.
The fixed $Y_i$ already respect $D\times\CP^2$, its complement,
and the finite flag; retain these same complexes and ordered diagonals.
By \cref{bd:relative,bd:union}, these give compatible controlled
quadratic pairs supported on the designated target sides. Their
interface is $\partial D\times\CP^2$; it is separated from $Z_0$.
By \cref{lem:union-boundary}, the first boundary is this interface,
and each subsequent boundary is its next equator times $\CP^2$.
After $n$ boundaries, projecting onto one component retains exactly
one copy of $\CP^2$, with sign $\varepsilon$ independent of the index.
At point control, subdivision identifies it with the constant model.
The complex orientation gives middle form $\langle1\rangle$;
its half-unit quadratic lift has the same symmetrized form and
signature $1$. This proves \eqref{bd:CP2}.
\end{proof}

\section{The contradiction}\label{sec:contradiction}

\begin{proposition}[Signature-one realization]\label{prop:signature-one}
An effective continuous $\Zp$-action on a connected positive-dimensional
manifold without boundary gives a fixed compact free $C_p$-space $T$,
finite cyclic groups $G_i$ with index-$p$ subgroups $P_i$, and
$\alpha\in L_4(\mathcal A_{\mathbf G}(T))$ with
\begin{equation}\label{eq:signature-one}
 \sigma_{\mathrm{tail}}(\alpha)=[(1,1,\ldots)].
\end{equation}
\end{proposition}
\begin{proof}
Use the fixed geometric data and the compensated normalized corner
$\mathfrak d$. Naturality of the iterated boundary and the scalar
comparison give
\[
 U_*\Delta_T\mathfrak d=\Delta_{\mathbf1}\mathfrak m_f
   =\varepsilon[C_*(\CP^2;\Q),\phi_{\CP^2}/2],
 \qquad \varepsilon\in\{1,-1\}.
\]
Set $\alpha=\varepsilon\Delta_T\mathfrak d$.
The coefficient-of-identity trace identifies the ordinary signature
of each equivariant form with that of its scalar image, so its tail
signature is the constant sequence $1$.
\end{proof}

\begin{proof}[Proof of \cref{thm:main}]
A connected zero-dimensional manifold is a point. For a connected
manifold with boundary, its interior is dense and connected:
each boundary half-ball has connected interior and meets just one
interior component, so different components would separate the manifold.
Homeomorphisms preserve the interior by local homology, and fixing it
fixes its closure.
Restriction to the interior therefore preserves effectivity.

If the boundaryless action were effective, \cref{prop:signature-one}
would give a class whose signatures are eventually $1$.
\Cref{alg:divisibility} makes the same signatures eventually divisible
by $p$. A common tail gives $1\in p\Z$, a contradiction.
\end{proof}

\begin{proof}[Proof of \cref{cor:HS}]
The point case is immediate; otherwise restrict to the interior, and
write $G$ for the acting group. Let $K$ be a compact subgroup of $G$. The
invariant-neighborhood construction in \S\ref{geo:fixed-data} gives a
nonempty $K$-invariant connected second-countable open submanifold $U$.
Put $N_U=\ker(K\curvearrowright U)$. The compact quotient $K/N_U$ embeds
in the metrizable group $\operatorname{Homeo}(U)$, so it is second
countable. The standard compact $p$-adic reduction \cite{Pa19} and
\cref{thm:main} make $K/N_U$ a Lie group.
If $U\subset V$ is another such invariant submanifold, the compact Lie
group $N_U/N_V\subset K/N_V$ fixes $U$, so Newman's theorem on $V$
\cite[Theorem (Newman)]{Pa19} gives $N_U=N_V$.
Every point lies in such a $V$ containing $U$, by the same construction.
Thus $N_U$ fixes the whole manifold and is trivial. Hence every compact
subgroup of $G$ is Lie.

Gleason--Yamabe gives an open subgroup $G_0\subset G$ with a compact
normal subgroup $K$ and Lie quotient $G_0/K$ \cite{Gle52,Yam53}.
Both $K$ and $G_0/K$ have no small subgroups (NSS). Intersect the inverse
image of an NSS neighborhood in $G_0/K$ with an identity neighborhood
whose intersection with $K$ contains no nontrivial subgroup. The
intersection contains no nontrivial subgroup, so $G_0$, and hence $G$,
has no small subgroups. Therefore $G$ is Lie with its original topology
\cite{Gle52,Yam53}.
\end{proof}

\section*{AI disclosure}
Various large language models (LLMs), including from Anthropic, OpenAI,
and open-source projects, were used under human guidance at every stage and for nearly
all tasks, from understanding, planning, and ideation to research, preparation, and writing.
A custom LLM orchestration framework was used for research. Critical construction steps
were developed jointly by Opus 5.5 and GPT 6.1. A formalization of \cref{thm:main,cor:HS} under second countability assumptions,
completed in Lean~4 by Opus 5.5 and verified with Comparator and SafeVerify,
is available at \url{https://github.com/adbrw/HS_proof}. This work was carried out as an independent research project.

\begin{thebibliography}{Ran80II}

\bibitem[BS01]{BS01}
P. Balmer and M. Schlichting,
\emph{Idempotent completion of triangulated categories},
J. Algebra \textbf{236} (2001), no.~2, 819--834.
\doi{10.1006/jabr.2000.8529}.

\bibitem[CP95]{CP95}
G. Carlsson and E. K. Pedersen,
\emph{Controlled algebra and the Novikov conjectures for $K$- and $L$-theory},
Topology \textbf{34} (1995), no.~3, 731--758.
\doi{10.1016/0040-9383(94)00033-H}.

\bibitem[Gle52]{Gle52}
A. M. Gleason,
\emph{Groups without small subgroups},
Ann. of Math. (2) \textbf{56} (1952), no.~2, 193--212.
\doi{10.2307/1969795}.

\bibitem[Hat02]{Hat02}
A. Hatcher,
\emph{Algebraic Topology},
Cambridge University Press, Cambridge, 2002.

\bibitem[Mar99]{Mar99}
G. J. Martin,
\emph{The Hilbert--Smith conjecture for quasiconformal actions},
Electron. Res. Announc. Amer. Math. Soc. \textbf{5} (1999), 66--70.
\doi{10.1090/S1079-6762-99-00062-1}.

\bibitem[Pa13]{Pa13}
J. Pardon,
\emph{The Hilbert--Smith conjecture for three-manifolds},
J. Amer. Math. Soc. \textbf{26} (2013), no.~3, 879--899.
\doi{10.1090/S0894-0347-2013-00766-3}.

\bibitem[Pa19]{Pa19}
J. Pardon,
\emph{Totally disconnected groups (not) acting on two-manifolds},
in \emph{Breadth in Contemporary Topology}, D. T. Gay and W. Wu (eds.),
Proc. Sympos. Pure Math. \textbf{102}, American Mathematical Society,
Providence, RI, 2019, 187--193.
\doi{10.1090/pspum/102/01809}.

\bibitem[Ran80I]{Ran80I}
A. A. Ranicki,
\emph{The algebraic theory of surgery. I. Foundations},
Proc. London Math. Soc. (3) \textbf{40} (1980), no.~1, 87--192.
\doi{10.1112/plms/s3-40.1.87}.

\bibitem[Ran80II]{Ran80II}
A. A. Ranicki,
\emph{The algebraic theory of surgery. II. Applications to topology},
Proc. London Math. Soc. (3) \textbf{40} (1980), no.~2, 193--283.
\doi{10.1112/plms/s3-40.2.193}.

\bibitem[Ran89]{Ran89}
A. A. Ranicki,
\emph{Additive $L$-theory},
$K$-Theory \textbf{3} (1989), no.~2, 163--195.
\doi{10.1007/BF00533377}.

\bibitem[Ran92]{Ran92}
A. A. Ranicki,
\emph{Lower $K$- and $L$-Theory},
London Math. Soc. Lecture Note Ser. \textbf{178},
Cambridge University Press, Cambridge, 1992.

\bibitem[Ran99]{Ran99}
A. A. Ranicki,
\emph{Singularities, double points, controlled topology and chain duality},
Doc. Math. \textbf{4} (1999), 1--59.
\doi{10.4171/DM/52}.

\bibitem[Rob75]{Rob75}
J. W. Roberts,
\emph{Pointwise finite families of mappings},
Canad. Math. Bull. \textbf{18} (1975), no.~5, 767--768.
\doi{10.4153/CMB-1975-134-x}.

\bibitem[RS97]{RS97}
D. Repov\v{s} and E. V. \v{S}\v{c}epin,
\emph{A proof of the Hilbert--Smith conjecture for actions by Lipschitz maps},
Math. Ann. \textbf{308} (1997), no.~2, 361--364.
\doi{10.1007/s002080050080}.

\bibitem[Wei13]{Wei13}
C. A. Weibel,
\emph{The $K$-Book: An Introduction to Algebraic $K$-Theory},
Grad. Stud. Math. \textbf{145}, American Mathematical Society,
Providence, RI, 2013.

\bibitem[Yam53]{Yam53}
H. Yamabe,
\emph{A generalization of a theorem of Gleason},
Ann. of Math. (2) \textbf{58} (1953), no.~2, 351--365.
\doi{10.2307/1969792}.

% Keep the author address with the final bibliography entries.
\enlargethispage{2\baselineskip}
\end{thebibliography}
\end{document}
